Considere $n\in\mathbb{N}$, $n\geq 3$, então lembremos
\begin{align*}
  S_n=\{\sigma:\{1,\cdots,n\}\to\{1,\cdots,n\}:\sigma\text{ é bijetiva.}\}
\end{align*}
Sabemos que $|S_n|=n!$.\\
Então, dado  $\sigma\in S_n$ com $\sigma\neq 1$, podemos considerar a estructura cíclica de $\sigma$.
\begin{definition}
  Chamamos a $\sigma$ um cíclo de comprimendo $m$ se $\sigma$ é uma permutação do tipo
  \begin{align*}
    \sigma:G&\to G\\
    i_1&\to i_2,\\
    i_2&\to i_3,\\
    &\vdots\\
    i_m&\to i_1.
  \end{align*}
  Escrevemos $\sigma=(i_1,i_2,\cdots,i_m)$ ou $m$-cíclo. 
\end{definition}
Note que toda permutação $\sigma\neq 1$ em $S_n$ pode ser escrita como um produto de cíclos disjuntos (cíclos que não mvem elementos em común).\\
Dois cíclos disjuntos comutam entre sí, então ela é a estructura cíclica de $\sigma$, i.e, é a manera única a mas de ordem do $\sigma$ onde $\sigma$ se escreve como produto de cíclos disjuntos.
Exemplos: 
\begin{align*}
  (9357)=(3579)=(5793)=(7935)\\
  (3\,5\,7\,9)(2\,4)(8\, 10\, 11)=(3\, 5\, 7\, 9)(9\, 10\, 11)(2\, 4)
\end{align*}
Você pode fazer e seguinte exercicio.
\begin{align*}
  (359)(89)(34)=(34598).
\end{align*}
\begin{note}
  Todo $m$-cíclo pode ser escrito como um produto de [transposições] (cíclos de comprimendo $2$) não necesariamente disjuntos, por exemplo
  \begin{align*}
    (2359)=(29)(25)(23).
  \end{align*}
  Geralmente 
  \begin{align*}
    (i_1\, i_2\, \cdots\, i_m)=(i_1i_m)(i_1i_{m-1})\cdots (i_1i_2).
  \end{align*}
\end{note}
\begin{note}[Propriedades de $S_n$]\label{note:propriedades-gerador-Sn}
  Um conjunto gerador de $S_n$
  \begin{align*}
    \left\langle (ij):1\leq i\neq j\leq n \right\rangle=S_n.
  \end{align*}
  Temos as seguintes propriedades
  \begin{itemize}
    \item $S_n=\left\langle (12),(23),\cdots,((n-1) n) \right\rangle$.
    \item $S_n=\left\langle (12),(12\cdots n) \right\rangle$.
    \item Dada $\varsigma=(ij)$ transpoção com $i\leq j-1$, temos $(ij)=((j-1)j)(i(j-1))((j-1)j)$. 
  \end{itemize}
\end{note}
\begin{definition}[Permutações pares e impares]
  $\sigma\in S_n$ é \textbf{par} se pode ser escrita como um produto de uma cantidade par de transposições. Caso contrário, $\sigma$ é \textbf{impar}. 
\end{definition}
\begin{definition}[$A_n$]
  Definimos $A_n$ como
  \begin{align*}
    A_n:=\{\sigma\in S_n:\sigma\text{ é par.}\}
  \end{align*}
  Ele é um subgrupo de $S_n$ dito subgrupo \textbf{alternado de grau $n$} de $S_n$.\\
  Note que $|A_n|=\frac{n!}{2}$.
\end{definition}
\begin{note}[$A_n\normal S_n$]
  Note que $A_n\normal S_n$ para todo $n\geq 5$.\\
  Seja $\varsigma=(ij)$ uma transposição de $S_n$ defina
  \begin{align*}
    f:S_n&\to A_n\\
    \sigma&\to f(\sigma)= 
    \begin{cases}
      \sigma, &\text{ se } \sigma\text{ e par} \text{,} \\
      \varsigma\sigma, &\text{ se } \sigma\text{ e impar} .
    \end{cases}
  \end{align*}
  Logo $f$ é um homomorfismoo sobrejetivo, $\ker f=:\{1,\varsigma\}=\left\langle \varsigma \right\rangle$. Assim, $|S_n/\ker f|=2\cong \Img f$, logo $[S_n:A_n]=2$ e portanto $A_n\normal S_n$. 
\end{note}
Consequentemente $S_n$ não é um grupo simple, mas, vamos mostrar que $A_n$ é um grupo sumple com $n\geq 3$ e $n\neq 4$.\\ 
Em $S_4$ nos temos que $|A_4|=12$, logo $\tilde{K_4}=\{1,(12)(34),(13)(23),(14)(23)\}$ e $\tilde{K_4}\normal A_4$ (mais do fue isto $\tilde{K_4}\normal S_4$).
\begin{proposition}
  Duas permutações $\alpha,\beta\in S_n$ das conjugada se, e somente se, $\alpha$ e $\beta$ tem a mesma estructura cíclica.
\end{proposition}
\begin{proof} 
  Suponha que existe $\gamma\in S_n$ tal que $\alpha=\beta^{\sigma}$, vamos admitir que $\alpha$ é um $m$-cíclo, digamos $\alpha=(i_1i_2\cdots i_m)$, temos
  \begin{align*}
    \beta\gamma(i_1)&=\gamma\alpha(i_1)\\
    &=\gamma(i_2).\\
    \beta\gamma(i_2)&=\gamma\alpha(i_2)\\
    &=\gamma(i_3).\\
    &\vdots
  \end{align*}
  Concluimos que $\beta$ é o cíclo $(\gamma(i_1),\gamma(i_2)\cdots\gamma(i_m))$.

  Reciprocamente, sejan
  \begin{align*}
    \alpha&=(i_1^{(1)}i_2^{(1)}\cdots i_{k_1}^{(1)})\cdots(i_1^{(s)}i_{2}^{(s)}\cdots i_{k_s}^{(s)}),\text{ disjunto }\\
    \beta&=(j_1^{(1)}j_2^{(1)}\cdots j_{k_1}^{(1)})\cdots(j_1^{(s)}j_{2}^{(s)}\cdots j_{k_s}^{(s)}),\text{ disjunto }
  \end{align*}
  Defina $\gamma: i_t^{l}\to j_t^{l}$.
\end{proof}
\begin{theorem}[$A_n$ é simple]
  Seja $G=A_n$. Se $n\geq 5$ temos que $A_n$ é simple. 
\end{theorem}
\begin{proof} 
  Seja $\{1\}\neq N\normal A_n$, para mostrar que $N=A_n$ basta mostrar que $N$ contem um $3$-cíclo, ja que os $3$-cíclos são conjugação de outro $3$-cíclo.\\
  Considere $p$ primo tal que $p\mid |N|$.\\
  Pelo teorema de Cauchy, existe $\sigma\in N$ tal que $o(\sigma)=p$.\\
  Escrevemos $\sigma$ como produto de ciclos disjuntos $\varsigma_i$, $1\leq i\leq m$, logo $\sigma=\varsigma_1\cdots\varsigma_m$.\\
  Deta forma, cada $\varphi_i$ é um $p$-cíclo.\\
  Analisamos os valorespossivels para $p$.
  \begin{enumerate}
    \item $p=2$.\\
      uma transposiçõe e assim, $m\geq 2$ par, escreva
      \begin{align*}
        \sigma=(ab)(cd)\varsigma_3\cdots\varsigma_m\in N.
      \end{align*}
      Considere $\gamma=(abc)$, temo que como $N$ é normal se pode calcular
      \begin{align*}
        \sigma^{\gamma}\sigma^{-1}&=(ab)^{\gamma}(cd)^{\gamma}\varsigma_3^{\gamma}\cdots\varsigma_{m}^{\gamma}\cdot \sigma^{-1}\\
        &=(ab)^{\gamma}(cd)^{\gamma}(ab)(cd)\\
        &=(ac)(bd)\in N.
      \end{align*}
      Agora, como $n\geq 5$, podemos tomar $k\notin\{a,b,c,d\}$, consideremos $\theta=(ack)$ e $\alpha=(ac)(bd)$ e temos que $(akc)=\alpha^{\theta}\alpha\in N$.
    \item $p=3$.\\
      Neste caso, cada $\varsigma_i$ é um $3$-cíclo, se $m=1$, acabou. Caso contrário, se $m\geq 2$, escreva $\sigma=(abc)(def)\varsigma_3\cdots\varsigma_m$. Considere o $3$-cíclo $\gamma=(bcd)$ e, como $N$ é normal se cumpre que $\sigma^{\gamma}\sigma^{-1}\in N$ e é um $5$-cíclo, este caso será resolvido a partir de na próxima situação.
    \item $p>3$.\\
      Consider $\varsigma_1=(a_1a_2\cdots a_p)$: $p$-cíclo em $N$ e $\gamma=(a_2a_3a_4)$ (observe que $\gamma$ comuta com todos $\varsigma_i$ com $i\geq 2$). Temos que
      \begin{align*}
        \sigma^{\gamma}\sigma^{-1}=(a_2a_3a_5)\in N.
      \end{align*}
      Isto mostra o resultado.
  \end{enumerate}
  Em conclução $N=A_n$.
\end{proof}
